% luamathaccents-doc.tex
% Documentation for the luamathaccents package.
% Copyright 2026 Joris Pinkse
% Distributed under the LaTeX Project Public License v1.3c or later.
% Compile with: lualatex luamathaccents-doc.tex
\documentclass{article}
\usepackage{geometry}
\usepackage{unicode-math}
% STIX Two Math and JuliaMono cover all the combining characters and
% accent glyphs shown in this documentation (the defaults, Latin Modern
% Math and Latin Modern Mono, miss a few of the rarer ones).  Both ship
% with TeX Live; if unavailable, compilation still succeeds but some
% glyphs come out blank.
\IfFontExistsTF{STIXTwoMath-Regular.otf}
  {\setmathfont{STIXTwoMath-Regular.otf}}{}
\IfFontExistsTF{JuliaMono-Regular.ttf}
  {\setmonofont{JuliaMono-Regular.ttf}%
     [Scale=MatchLowercase, BoldFont=JuliaMono-Bold.ttf]}{}
\usepackage{metalogo}
\usepackage{booktabs}
\usepackage{longtable}
\usepackage{hyperref}
\usepackage{luamathaccents}

\title{The \textsf{luamathaccents} package\\[1ex]
  \large Unicode combining accents in math input, v1.0}
\author{Joris Pinkse\\\texttt{pinkse@gmail.com}}
\date{2026/07/26}

\begin{document}
\maketitle

\section{Introduction}

Modern editors and input methods make it easy to type accented
mathematical symbols directly, using Unicode \emph{combining
characters}.  For instance, typing \texttt{x\textbackslash hat}
followed by \texttt{Tab} in VS~Code (with the Julia or Unicode Latex
extension) shows an $x$ with a hat on top.  However, \LaTeX{} does not
interpret this combination correctly.

What this package does is to make \LaTeX{} interpret an accented
character in the source correctly.  This package requires \LuaLaTeX{}.

\section{Usage}

Load the package in the preamble of a document compiled with
\texttt{lualatex}:

\begin{verbatim}
\usepackage{luamathaccents}
\end{verbatim}

The rewriting is active from that point on.  Two commands control it:

\begin{description}
  \item[\texttt{\textbackslash LuaMathAccentsOff}] stop rewriting input
    lines (takes effect from the next input line);
  \item[\texttt{\textbackslash LuaMathAccentsOn}] resume rewriting.
\end{description}

Most of the generated commands are the classic math accents
(\verb|\hat|, \verb|\tilde|, \verb|\vec|, \ldots).  A number of the
rarer ones (\verb|\ovhook|, \verb|\candra|, \verb|\annuity|, \ldots)
are provided by the \textsf{unicode-math} package, which you should
load if you use those characters.  Since you are already compiling
with \LuaTeX{}, loading \textsf{unicode-math} is recommended anyway.

\section{Example}

\LuaMathAccentsOff
The input
\begin{verbatim}
$x̂ + ỹ = z̄$, force $F⃗ = m a⃗$, and average $v̄$.
\end{verbatim}
\LuaMathAccentsOn
is typeset as
\begin{quote}
$x̂ + ỹ = z̄$, force $F⃗ = m a⃗$, and average $v̄$.
\end{quote}

\section{Supported combining characters}

\begin{longtable}{llll}
\toprule
Code & Unicode name (\textsc{combining} \ldots) & Command & Result\\
\midrule
\endhead
\texttt{U+0300} & grave accent & \verb|\grave| & $\grave{x}$\\
\texttt{U+0301} & acute accent & \verb|\acute| & $\acute{x}$\\
\texttt{U+0302} & circumflex accent & \verb|\hat| & $\hat{x}$\\
\texttt{U+0303} & tilde & \verb|\tilde| & $\tilde{x}$\\
\texttt{U+0304} & macron & \verb|\bar| & $\bar{x}$\\
\texttt{U+0305} & overline & \verb|\bar| & $\bar{x}$\\
\texttt{U+0306} & breve & \verb|\breve| & $\breve{x}$\\
\texttt{U+0307} & dot above & \verb|\dot| & $\dot{x}$\\
\texttt{U+0308} & diaeresis & \verb|\ddot| & $\ddot{x}$\\
\texttt{U+0309} & hook above & \verb|\ovhook|\textsuperscript{*} & $\ovhook{x}$\\
\texttt{U+030A} & ring above & \verb|\ocirc|\textsuperscript{*} & $\ocirc{x}$\\
\texttt{U+030C} & caron & \verb|\check| & $\check{x}$\\
\texttt{U+0310} & candrabindu & \verb|\candra|\textsuperscript{*} & $\candra{x}$\\
\texttt{U+0312} & turned comma above & \verb|\oturnedcomma|\textsuperscript{*} & $\oturnedcomma{x}$\\
\texttt{U+0315} & comma above right & \verb|\ocommatopright|\textsuperscript{*} & $\ocommatopright{x}$\\
\texttt{U+031A} & left angle above & \verb|\droang|\textsuperscript{*} & $\droang{x}$\\
\texttt{U+0331} & macron below & \verb|\underline| & $\underline{x}$\\
\texttt{U+20D0} & left harpoon above & \verb|\leftharpoonaccent|\textsuperscript{*} & $\leftharpoonaccent{x}$\\
\texttt{U+20D1} & right harpoon above & \verb|\rightharpoonaccent|\textsuperscript{*} & $\rightharpoonaccent{x}$\\
\texttt{U+20D6} & left arrow above & \verb|\overleftarrow| & $\overleftarrow{x}$\\
\texttt{U+20D7} & right arrow above & \verb|\vec| & $\vec{x}$\\
\texttt{U+20DB} & three dots above & \verb|\dddot|\textsuperscript{*} & $\dddot{x}$\\
\texttt{U+20E1} & left right arrow above & \verb|\overleftrightarrow|\textsuperscript{*} & $\overleftrightarrow{x}$\\
\texttt{U+20E7} & annuity symbol & \verb|\annuity|\textsuperscript{*} & $\annuity{x}$\\
\texttt{U+20E9} & wide bridge above & \verb|\widebridgeabove|\textsuperscript{*} & $\widebridgeabove{x}$\\
\texttt{U+20F0} & asterisk above & \verb|\asteraccent|\textsuperscript{*} & $\asteraccent{x}$\\
\bottomrule
\end{longtable}

\noindent
Commands marked \textsuperscript{*} are provided by
\textsf{unicode-math} (or, for \verb|\dddot|, also by
\textsf{amsmath}).

\section{Caveats}

\begin{itemize}
  \item The rewriting operates on \emph{every} input line, in text
    mode as well as math mode, and also inside \texttt{verbatim}
    material, since it happens before TeX tokenizes the line.  Use
    \verb|\LuaMathAccentsOff| around text where combining characters must
    be left alone.
  \item Only \emph{combining} sequences (Unicode normalization form
    NFD) are recognized.  Precomposed characters such as
    \texttt{U+0227} (\textsc{latin small letter a with dot above}) are
    a single codepoint and pass through unchanged.
  \item A combining character with no preceding base character on the
    same line is passed through unchanged.
  \item Only one accent per base character is converted; in a sequence
    of several combining characters, the ones after the first are
    passed through unchanged (whether they produce any output then
    depends on the font).
\end{itemize}

\section{License}

Copyright 2026 Joris Pinkse.  This work may be distributed and/or
modified under the conditions of the \LaTeX{} Project Public License,
either version~1.3c of this license or (at your option) any later
version (\url{https://www.latex-project.org/lppl.txt}).

\end{document}
